% council-handbook.tex
% =====================================================================
% A LaTeX recreation of "Council -- The Handbook (Review Edition)".
%
% This is the example document for gigvim's TeX/LaTeX support. It is a
% faithful, self-contained rebuild of the reviewed Council rulebook:
% dark title bands, cream "proposal" sidebars marked with a decision
% flag, small-caps tracked headers, pull-quotes and card grids -- the
% kind of layout that shows off what LaTeX can typeset, and gives
% texlab (LSP) and tex-fmt (formatter) real material to work on.
%
% Build (either works; tectonic needs no bibtex/latexmk dance):
%   tectonic council-handbook.tex
%   latexmk -pdf council-handbook.tex
%
% In gigvim: open the file and save -- tex-fmt reformats on write
% (formatOnSave is on globally); :LspInfo shows texlab attached.
%
% Compiles with pdflatex/tectonic. Fonts are stock Latin Modern for
% portability; swap in ebgaramond/libertinus for a closer match.
% =====================================================================

\documentclass[11pt]{article}

\usepackage[T1]{fontenc}
\usepackage[utf8]{inputenc}
\usepackage{lmodern}
\usepackage[a4paper,margin=22mm,bottom=24mm]{geometry}
\usepackage{microtype}          % \textls letterspacing, better justification
\usepackage[most]{tcolorbox}    % boxes, cards, rasters, skins
\usepackage{pifont}             % \ding dingbats (flag/star stand-ins)
\usepackage{amssymb}            % \circlearrowleft
\usepackage[normalem]{ulem}     % \sout / \uline for review marks
\usepackage{enumitem}
\usepackage{fancyhdr}
\usepackage{needspace}
\usepackage[hidelinks]{hyperref}

% ---------------------------------------------------------------------
% Palette (sampled from the reviewed edition)
% ---------------------------------------------------------------------
\definecolor{navy}{HTML}{16293B}
\definecolor{cream}{HTML}{ECE3CC}
\definecolor{gold}{HTML}{C39A46}
\definecolor{teal}{HTML}{4E7E86}
\definecolor{inkgreen}{HTML}{2E7D4F}
\definecolor{inkred}{HTML}{B23A3A}
\definecolor{stone}{HTML}{6B7480}

% ---------------------------------------------------------------------
% Glyph stand-ins (portable under pdflatex)
% ---------------------------------------------------------------------
\newcommand{\flag}{\textcolor{gold}{\ding{110}}}      % decision marker
\newcommand{\deco}{\textcolor{gold}{\ding{72}}}       % ornament
\newcommand{\reloop}{$\circlearrowleft$}

% Review marks: struck (removed) text and added text.
\newcommand{\cut}[1]{\textcolor{inkred}{\sout{#1}}}
\newcommand{\add}[1]{\textcolor{inkgreen}{\uline{#1}}}
\newcommand{\hlt}[1]{\colorbox{gold!22}{#1}}

% ---------------------------------------------------------------------
% Typographic helpers
% ---------------------------------------------------------------------
\newcommand{\eyebrow}[1]{{\sffamily\footnotesize\bfseries\textcolor{teal}%
    {\textls[140]{\MakeUppercase{#1}}}}}

% Section header, e.g. \sect{0.3 The Turn of the Tide --- The Core Loop}
\newcommand{\sect}[1]{%
  \par\addvspace{1.15\baselineskip}\needspace{4\baselineskip}%
  \noindent{\sffamily\bfseries\small\textcolor{gold}%
    {\textls[110]{\S\ \MakeUppercase{#1}}}}%
  \par\addvspace{0.35\baselineskip}\nobreak}

% Card head: small tracked label over a serif title
\newcommand{\cardhead}[2]{\leavevmode{\sffamily\scriptsize\textls[120]{\MakeUppercase{#1}}}\\[1mm]{\large #2}\par\vspace{0.6mm}}

% ---------------------------------------------------------------------
% Book title band (dark navy full-width)
% ---------------------------------------------------------------------
% \booktitle{eyebrow}{TITLE}{book the ...}{italic tagline}
\newcommand{\booktitle}[4]{%
  \begin{tcolorbox}[enhanced,sharp corners,boxrule=0pt,colback=navy,
      colframe=navy,coltext=cream,left=8mm,right=8mm,top=11mm,bottom=9mm,
      before skip=4mm,after skip=8mm]
    \centering
    {\sffamily\scriptsize\textls[240]{\MakeUppercase{#1}}}\\[4mm]
    {\fontsize{40}{44}\selectfont\textls[40]{#2}}\\[2.5mm]
    {\sffamily\footnotesize\textls[240]{\MakeUppercase{#3}}}\\[3.5mm]
    {\itshape\color{cream!88}#4}
  \end{tcolorbox}}

% ---------------------------------------------------------------------
% Pull-quote (teal left rule, italic)
% ---------------------------------------------------------------------
\newenvironment{pullquote}
  {\par\addvspace{0.7\baselineskip}%
   \begin{tcolorbox}[blanker,borderline west={2.5pt}{0pt}{teal},
       left=5mm,top=0.5mm,bottom=0.5mm,before skip=2mm,after skip=2mm]%
   \itshape\large}
  {\end{tcolorbox}\par\addvspace{0.5\baselineskip}}

% ---------------------------------------------------------------------
% Proposal sidebar (cream, gold left rule) + its header line
% ---------------------------------------------------------------------
\newtcolorbox{proposalbox}{enhanced,breakable,sharp corners,colback=cream,
  colframe=cream,coltext=black,boxrule=0pt,
  borderline west={2.5pt}{0pt}{gold},
  left=5mm,right=5mm,top=3.5mm,bottom=3.5mm,before skip=3.5mm,after skip=3.5mm}

\newcommand{\prophead}[1]{%
  \noindent{\sffamily\bfseries\footnotesize\textcolor{gold!75!black}%
    {\flag\ \textls[60]{\MakeUppercase{#1}}}}%
  \nobreak\hfill{\sffamily\itshape\scriptsize\textcolor{stone}{for your sign-off}}%
  \par\vspace{2mm}}

% Highlighted "aside" panel (a plain cream note without the flag)
\newtcolorbox{notebox}{enhanced,breakable,sharp corners,colback=cream,
  colframe=cream,coltext=black,boxrule=0pt,
  left=5mm,right=5mm,top=3.5mm,bottom=3.5mm,before skip=3.5mm,after skip=3.5mm}

% ---------------------------------------------------------------------
% Running foot
% ---------------------------------------------------------------------
\pagestyle{fancy}
\fancyhf{}
\renewcommand{\headrulewidth}{0pt}
\fancyfoot[C]{\sffamily\scriptsize\textcolor{stone}%
  {\textls[120]{COUNCIL · THE HANDBOOK · REVIEW EDITION}}}
\fancyfoot[R]{\sffamily\scriptsize\textcolor{stone}{\thepage}}

\setlength{\parskip}{0.55\baselineskip}
\setlength{\parindent}{0pt}

\begin{document}

% =====================================================================
% FRONT MATTER -- how to give notes
% =====================================================================
\thispagestyle{fancy}

{\sffamily\scriptsize\textls[220]{\MakeUppercase{Council · Review Edition}}}\par
\vspace{1mm}
{\fontsize{34}{38}\selectfont The Handbook}\par
\vspace{1mm}
{\itshape\large\color{stone}Proposed revisions \& open proposals --- marked for your review}\par
\vspace{1mm}
{\color{gold}\rule{34mm}{1.2pt}}\par
\vspace{3mm}

\begin{tcolorbox}[enhanced,sharp corners,colback=navy,colframe=navy,
    coltext=cream,boxrule=0pt,borderline west={2.5pt}{0pt}{gold},
    left=5mm,right=5mm,top=4mm,bottom=4mm]
  {\sffamily\bfseries\footnotesize\textcolor{gold}%
    {\textls[120]{HOW TO GIVE ME NOTES YOU'VE HANDWRITTEN}}}\par\vspace{2mm}
  This whole document was rebuilt from notes you wrote by hand. Here's how to
  make the next batch read cleanly:
  \begin{enumerate}[leftmargin=7mm,itemsep=1.2mm,topsep=1.5mm]
    \item \textbf{Use pen, and use colour.} Blue, red, green --- any non-black
      ink. Colour lets me lift your marks cleanly off the black print. Pencil
      works but is fainter; a coloured pen is best.
    \item \textbf{Point before you talk.} Circle or box the exact word/line you
      mean, then draw a short line out to your note in the margin. That tells me
      \emph{what} the note is about, not just \emph{that} there is one.
    \item \textbf{One mark, one meaning.} Highlight = ``this matters.''
      Strike-through = ``cut this.'' Caret (\textasciicircum) with a word =
      ``insert here.'' A question in the margin = ``decide this with me.''
    \item \textbf{Export as PDF, not a photo.} Annotating and exporting a PDF
      keeps your ink as crisp vector strokes --- those extract perfectly. A
      phone photo of paper works too, just less sharp.
    \item \textbf{Big decisions in words.} For the \flag\ proposals, a margin
      ``yes,'' ``no,'' or ``option (a)'' is enough --- or dictate them to me
      like you did this round, which worked great.
  \end{enumerate}
  \vspace{1mm}
  {\itshape\color{cream!85}Then just upload the marked-up PDF and say ``read my
  notes.'' I'll pull every mark, apply the settled ones, and bring the open
  ones back to you.}
\end{tcolorbox}

\vspace{2mm}
{\sffamily\scriptsize\textcolor{stone}{\textls[120]{TAP THE CONTENTS OR ANY \S\ TO JUMP · \flag\ MARKS A DECISION FOR YOU}}}

% =====================================================================
% COVER
% =====================================================================
\clearpage
\booktitle{A roleplaying game · played with words, not dice}{COUNCIL}%
  {A Player's Handbook · Printed Edition · One Voyage, One Volume}%
  {strive and pay, or be given --- but keep hope, whatever the weather}

\vspace{4mm}
\begin{center}\deco\end{center}

% =====================================================================
% CONTENTS
% =====================================================================
\clearpage
\eyebrow{\deco\ Contents \deco}\par
\vspace{1mm}
{\fontsize{26}{30}\selectfont Where to find things}\par
\vspace{3mm}

{\large How to Play}\quad{\sffamily\scriptsize\textcolor{stone}{\textls[120]{FIVE MINUTES · READ THIS FIRST}}}

\newcommand{\bookrow}[3]{%
  \vspace{2mm}{\large\textcolor{gold}{#1}\ \ #2}\par
  {\sffamily\footnotesize\textcolor{stone}{#3}}\par}

\bookrow{I}{Foundations}{\S0.1 The Thesis \quad \S0.2 Who Sails Under It \quad \S0.3 The Core Loop \quad \S0.4 The Eight Timbers \quad \S0.5 Keel \& Rigging}
\bookrow{II}{The Resolution Core}{\S2.1 The Two Answers \quad \S2.2 Grace \quad \S2.3 Price \quad \S2.4 The Seals \quad \S2.5 The Roads \quad \S2.6 The Reward of Cleverness}
\bookrow{III}{Works \& the Tide}{\S3.1 Declaring a Work \quad \S3.2 The Tide \quad \S3.3 The Ledger \quad \S3.4 Squalls \quad \S3.5 Forcing a Work \quad \S3.6 Completion \& Work Orders \quad \S3.7 The Divine Work}
\bookrow{IV}{Governance}{\S4.1 The Council \quad \S4.2 Who Sits \quad \S4.3 Deputies \quad \S4.4 Interrupts \quad \S4.5 The Override \quad \S4.6 Absence Beyond Reach}
\bookrow{V}{Teach \& Ship}{\S5.1 Session Zero \quad \S5.2 Writing Seals \quad \S5.3 Giving Grace \quad \S5.4 The Signed Squall \quad \S5.5 Hope \& Harm \quad \S5.6 What to Print}

% =====================================================================
% PROPOSALS INDEX + HOW TO PLAY
% =====================================================================
\clearpage
\begin{proposalbox}
  \prophead{Proposals awaiting sign-off}
  {\sffamily\footnotesize\textcolor{teal}{%
    1 · Turning the watch \quad 2 · Ledger meat \quad 3 · Landis' gifts \quad
    4 · Sealed favor \quad 5 · Between-session play \quad 6 · Trigger Points \quad
    7 · Three tiers \quad 8 · Dependent Works \quad 9 · The Fasces \quad
    10 · The Minute \quad 11 · Reversion \quad 12 · Absence rewrite \quad
    13 · The creed \quad 14 · Clue sourcing}}
\end{proposalbox}

\begin{tcolorbox}[enhanced,sharp corners,colback=navy,colframe=navy,
    coltext=cream,boxrule=0pt,left=6mm,right=6mm,top=6mm,bottom=6mm]
  \centering
  {\sffamily\scriptsize\textls[200]{HOW TO PLAY · A GUIDE FOR THE CREW · FIVE MINUTES}}\\[2mm]
  {\fontsize{30}{34}\selectfont Before You Sit Down}
\end{tcolorbox}

\begin{pullquote}
  There are no dice here, and no character sheets. There is only the story we
  tell together --- and a handful of promises about how it answers back.
\end{pullquote}

\begin{enumerate}[leftmargin=8mm,itemsep=1.5mm]
  \item \emph{You live a scene.} You speak and act as your character, and the
    keeper plays the world around you. Most of the time, that is all there is
    --- and it is enough.
  \item \emph{A need arises the table can't just wish away.} A locked door. A
    hidden truth. A thing that must be built or found. Under dice you'd roll.
    Here, you don't guess the magic question --- you simply \emph{reach} for the
    answer, and the world provides it by one of two ways below.
  \item \emph{What comes back changes the next scene.} And the tale sails on.
    That wheel --- scene, answer, scene --- is the whole of the game.
\end{enumerate}

\begin{tcolorbox}[blanker,before skip=3mm,after skip=3mm]
\begin{tcbraster}[raster columns=2,raster equal height=rows,
    raster column skip=4mm,sharp corners,boxrule=0pt,
    colback=navy,colframe=gold,coltext=cream,toprule=1.4pt,
    left=3.5mm,right=3.5mm,top=3mm,bottom=3mm]
  \tcbitem \cardhead{}{By Price} What you \emph{pay for}. You may seize an
    answer by spending something of your own --- a truth confessed, a new
    trouble taken on, a little time lost. Price is yours to demand. It always
    costs the self, and it is always true.
  \tcbitem \cardhead{}{By Grace} What you are \emph{given}. The keeper may hand
    you an answer freely --- unearned, unasked, because you are stuck and the
    story loves you. Grace cannot be bought or demanded. It falls where it
    falls, and it falls most on those who keep trying and keep hoping.
\end{tcbraster}
\end{tcolorbox}

\begin{center}\itshape Strive and pay, or be given. Most tales are told in the
weather between the two.\end{center}

% ---- You cannot be truly stuck ----
\clearpage
\begin{notebox}
  {\sffamily\bfseries\footnotesize\textcolor{gold!75!black}%
    {\textls[100]{YOU CANNOT BE TRULY STUCK}}}\par\vspace{1.5mm}
  Behind every real obstacle the keeper has already written three clues --- a
  gentle \emph{nudge}, a plain \emph{bearing}, and, at the last, the
  \emph{answer} itself. They are never improvised and never withheld forever.
  However you reach for help --- by cleverness, by price, by grace --- one of
  these comes. No riddle in Council is a locked room with no key. That is a
  promise.
\end{notebox}

\sect{The One Thing Worth Keeping}
Council is built to reward \emph{hope}. The world bends kinder to a crew that
believes it can be made better, and grace runs toward them like water downhill.
The game can wound you --- it can cost you dearly, and grieve you truly --- but
it only truly breaks the one who has stopped hoping. Keep hope, and you are
never past saving.

\begin{pullquote}Despair is the only door that locks from the inside.\end{pullquote}

\sect{One Stranger You Should Know}
There is a madman named \emph{Landis}, who is real only in water. Find still
water in any world --- a sea, a cistern, a cup --- look in, and ask him with
grace, and he will sell you a true answer at a steep and personal price. He is
not the main road, and you will not need him often. But he is always there, in
every water, and every answer he gives is true. More of him waits in the pages
that follow.

% =====================================================================
% BOOK I -- FOUNDATIONS
% =====================================================================
\clearpage
\booktitle{Council · A Player's Handbook}{FOUNDATIONS}{Book the First}%
  {being the whole of the game, set down plain, that any hand aboard may read it}

\sect{0.1 The Thesis}
\begin{pullquote}Council is for tables who left the dice behind --- and found
they had left something else behind with them.\end{pullquote}
When a crew stops rolling and starts telling, the story runs warmer and the
arithmetic goes quiet --- and then, one night, a locked door appears that no one
at the table happens to be clever enough to open. Under dice, a character
\emph{rolled to know things their player did not.} Take the dice away and that
quiet gift goes with them: now a player can learn only what they personally
thought to ask, and the keeper is left dangling hints, praying someone guesses
the shape of the answer.

\textbf{Council exists to hand that gift back without handing the dice back.} It
is a game of insight, consequence, and shared command that runs entirely on
words. One question steers every rule in this handbook: \emph{how does a soul
come to know a thing without guessing the exact right way to ask?} All that
follows is an answer.

\begin{pullquote}You will not roll to see if you succeed. You will decide what
you are willing to pay to find out. --- the first principle\end{pullquote}

\sect{0.2 Who Sails Under It}
\begin{tcbraster}[raster columns=2,raster equal height=rows,raster column skip=4mm,
    sharp corners,boxrule=0pt,colback=cream,colframe=cream,coltext=black,
    left=3.5mm,right=3.5mm,top=3mm,bottom=3mm]
  \tcbitem \emph{Crews who have already drifted.} Long campaigns that quietly
    abandoned sheets and initiative and now live on narration and trust. Council
    names what they are already doing and gives it teeth.
  \tcbitem \emph{Keepers weary of dangling hints.} Game masters who love the
    shared story but keep striking the reef where a puzzle stalls because no one
    has yet said the magic thing.
\end{tcbraster}

\sect{0.3 The Turn of the Tide --- The Core Loop}
Play turns through three stations. A \emph{Scene} raises a need the table cannot
simply declare true. The world answers by one of two roads --- \emph{the Water}
(buy insight now, at a price) or \emph{the Tide} (commit the work, let time
carry it) --- and whatever comes back walks into the next Scene, changed. That
is the entire machine. All else is rigging hung upon it.

\begin{tcbraster}[raster columns=3,raster equal height=rows,raster column skip=3mm,
    sharp corners,boxrule=0pt,colback=cream,colframe=cream,coltext=black,
    toprule=1.4pt,colupper=black,
    left=3.5mm,right=3.5mm,top=3mm,bottom=3mm]
  \tcbitem[colframe=gold] \cardhead{Station One}{The Scene} A need arises the
    table cannot just wish true.
  \tcbitem[colframe=gold] \cardhead{Road A}{The Water} A truth, a trouble, a
    scene, for a real answer.
  \tcbitem[colframe=teal] \cardhead{Road B}{The Tide} Commit a Work; the world
    moves it a watch at a time.
\end{tcbraster}
\begin{center}\small\textcolor{stone}{\reloop\ \ and whatever the world returns
walks back into the next Scene, changed\ \ \reloop}\end{center}

\sect{0.4 The Eight Timbers}
The fixed vocabulary of the game --- the timbers of the hull. These words mean
the same at every Council table; a keeper may re-rig the world above, but not
rename the timbers below.

\begin{tcolorbox}[enhanced,sharp corners,colback=cream,colframe=cream,
    coltext=black,boxrule=0pt,left=4mm,right=4mm,top=3mm,bottom=3mm]
\begin{description}[leftmargin=13mm,style=nextline,itemsep=1.2mm,font=\normalfont]
  \item[\textcolor{gold}{\bfseries I}\ \ The Water] Landis, real only in water,
    who sells true answers at a price.
  \item[\textcolor{gold}{\bfseries II}\ \ The Council] Shared command with a
    final say; every vote and dissent minuted.
  \item[\textcolor{gold}{\bfseries III}\ \ Works] What a crew sets in motion
    between scenes: an aim, the hands, the stakes.
  \item[\textcolor{gold}{\bfseries IV}\ \ The Tide] Time itself. \cut{Every
    session advances every open Work one watch.} \add{The keeper turns the
    watch; a session may pass one, or several.}
  \item[\textcolor{gold}{\bfseries V}\ \ The Ledger] What a Work spends: Hands,
    Power, Standing. Nothing pinned twice.
  \item[\textcolor{gold}{\bfseries VI}\ \ Squalls] Sealed complications carried
    by the larger Works, breaking when the weave turns against you.
  \item[\textcolor{gold}{\bfseries VII}\ \ The Seals] Pre-written clues ---
    Nudge, Bearing, Answer --- sealed and shown, in order.
  \item[\textcolor{gold}{\bfseries VIII}\ \ Price \& Grace] The two ways an
    answer comes: seized with a piece of self, or given freely.
\end{description}
\end{tcolorbox}

\sect{0.5 The Keel \& the Rigging}
Council is a keel, not a world. A table brings its own fiction --- a drowned
empire, a generation ship, a haunted county --- and rigs it to the timbers
below. But one timber is not the table's to swap. \emph{Landis is fixed.} The
madman on the sea, real only in water, who sells true answers at a price, stands
in every Council game that has been or will be run.

\begin{tcbraster}[raster columns=2,raster equal height=rows,raster column skip=4mm,
    sharp corners,boxrule=0pt,left=3.5mm,right=3.5mm,top=3mm,bottom=3mm]
  \tcbitem[colback=navy,colframe=gold,coltext=cream,toprule=1.4pt]
    \cardhead{The Keel --- Fixed}{} Landis \& the Water. The eight timbers and
    their names. The turn of the tide. Price and grace. The Seals. Mercy.
  \tcbitem[colback=cream,colframe=teal,coltext=black,toprule=1.4pt]
    \cardhead{The Rigging --- Yours}{} The world, its peoples, what Power and
    Standing are made of, what the riddles guard.
\end{tcbraster}
\begin{pullquote}Change everything you like. Leave me the water. --- L.F.\end{pullquote}

% =====================================================================
% BOOK II -- THE RESOLUTION CORE
% =====================================================================
\clearpage
\booktitle{Council · A Player's Handbook}{THE RESOLUTION CORE}{Book the Second}%
  {how the game answers, when there is no die to throw}

\begin{pullquote}A die answered a question the table could not: what is true,
and what do you learn? Council keeps the question and throws away the die. In
its place stand two hands --- one that takes, and one that gives. Learn these
two and you have learned the game.\end{pullquote}

\sect{2.1 The Two Answers}
When a scene raises a need the table cannot simply declare true, the world
answers in one of two ways, and only two. An answer is either \emph{seized} ---
bought with something of your own --- or it is \emph{given}, freely, by the hand
that made this world. The first is \emph{price}. The second is \emph{grace}.

\begin{tcbraster}[raster columns=2,raster equal height=rows,raster column skip=4mm,
    sharp corners,boxrule=0pt,colback=navy,colframe=gold,coltext=cream,
    toprule=1.4pt,left=3.5mm,right=3.5mm,top=3mm,bottom=3mm]
  \tcbitem[colframe=gold] \cardhead{}{Price} You pay, and you may demand. It
    costs the self and is always true.
  \tcbitem[colframe=teal] \cardhead{}{Grace} You are given, and may not demand.
    It costs you nothing and cannot be earned.
\end{tcbraster}

\sect{2.2 Grace --- The Given Answer}
The keeper is the god of this small creation --- not a tyrant, but a maker, who
set its laws and largely lets them run. \emph{Grace is the maker's one free
hand:} the single thing in the game not bound by the world's own rules. An
answer given as grace is unearned and unpriced.

Grace \emph{cannot be bought, banked, or demanded.} A good-faith struggle
\emph{invites} it, and hope draws it near --- but nothing compels it. It may
fall on the crew who tried hardest; it may fall, just as well, on the one who
tried least and needed it most. It falls where it will.

Grace is \emph{not counted.} There is no pool, no token, no tally at the table's
edge. It has exactly one law that binds it: \emph{mercy}. If a crew is truly
aground, grace comes, because the maker will not let the story drown.

And grace comes \emph{quietly}. It is a nudge slid across the table,
unremarked. The keeper need not announce it, and mostly should not. The crew
will feel they were lucky. Let them.

\begin{pullquote}Grace is not the keeper being kind against the rules. It is the
one rule that kindness was allowed to keep.\end{pullquote}

\sect{2.3 Price --- The Seized Answer}
Where grace is given, price is \emph{taken} --- by you, from you. When you must
know a thing now and are willing to bleed for it, you may seize the answer and
pay in the only coin the game accepts: something of your own self.

The dearest altar of price is \emph{the Water} --- Landis, real only in water,
who trades true answers for pieces of the one who asks. But mark this well:
\emph{the Water is a rare road, not a daily one.} When you do go down to the
water, the price climbs with each answer of a single visit:

\begin{tcolorbox}[enhanced,sharp corners,colback=cream,colframe=cream,
    coltext=black,boxrule=0pt,left=4mm,right=4mm,top=3mm,bottom=3mm]
\begin{description}[leftmargin=18mm,style=nextline,itemsep=1.2mm,font=\sffamily\footnotesize\bfseries\color{gold!75!black}]
  \item[FIRST] \emph{A truth.} Confess a secret, or commit a new fact about
    yourself, aloud. Canon the moment it is spoken.
  \item[SECOND] \emph{A trouble.} The answer is true, and arrives lashed to a
    fresh complication.
  \item[THIRD] \emph{The madness bleeds.} You glimpse a thing the keeper knows
    and you should not, and cannot say why.
  \item[ALWAYS] \emph{A scene.} To sit with Landis costs time; the weave moves
    on while you do.
\end{description}
\end{tcolorbox}

\sect{2.4 The Seals --- What Both Hands Reach For}
Price and grace do not invent their answers on the spot. \emph{Behind any
obstacle worth being stuck on lie three clues,} set face-down on the table, in
order.

\begin{tcbraster}[raster columns=3,raster equal height=rows,raster column skip=3mm,
    sharp corners,boxrule=0pt,colback=cream,colframe=navy,coltext=black,
    toprule=1.4pt,left=3.5mm,right=3.5mm,top=3mm,bottom=3mm]
  \tcbitem \cardhead{The First}{The Nudge} A single honest word or image.
  \tcbitem \cardhead{The Second}{The Bearing} What to do next --- never why.
  \tcbitem \cardhead{The Third}{The Answer} Plain as porridge.
\end{tcbraster}
They open \emph{in order}, one at a time. Price opens the next by paying for it.
Grace opens the next by giving it --- \hlt{most often the Nudge, for nothing.}

\sect{2.5 The Roads to a Clue}
A crew that means to pay has more than one road to the next seal --- the Water,
being dearest, is the one they should walk least:

\begin{tcolorbox}[enhanced,sharp corners,colback=cream,colframe=cream,
    coltext=black,boxrule=0pt,left=4mm,right=4mm,top=3mm,bottom=3mm]
\begin{description}[leftmargin=24mm,style=sameline,itemsep=1.4mm,font=\normalfont\bfseries]
  \item[Honest search] Dig, read, question, pry. Paid in effort and a scene's
    time. \add{--- call it wit by another name.}
  \item[A Work] Let the Tide carry it (Book III). Slow, priced only in watches.
  \item[A letter] Ask one who is absent but would know. \cut{Slow as the post.}
    \add{Slow as the post --- and those who carry it may read what it says,
    unless it is sealed.}
  \item[The Water] Landis, and his \hlt{climbing price}. Walk it least.
\end{description}
\end{tcolorbox}

\begin{tcolorbox}[enhanced,sharp corners,colback=navy,colframe=navy,
    coltext=cream,boxrule=0pt,left=4mm,right=4mm,top=3mm,bottom=3mm]
  \textbf{The rule that binds every road:} it opens \emph{the next Seal in
  order} --- never a later one, never two at once.
\end{tcolorbox}

\sect{2.6 The Reward of Cleverness}
Sometimes a crew closes out a whole labour --- a Work, a mission, an obstacle
worth being stuck on --- with its Seals never broken: no price paid, no grace
needed. Each Seal left closed is banked as \emph{Slack}: a due at the Water. The
next time the crew comes to Landis, that much of his usual price is simply
\emph{waived} --- one price forgiven for every Seal they earned --- and he
throws in something past the ordinary menu besides, a thing he judges cool
enough for the doing, never named in advance. Slack is spent only there, only
with him, and does not linger forever unspent.

\begin{pullquote}Solve it yourself, and the madman remembers you the next time
you're wet.\end{pullquote}

\begin{proposalbox}
  \prophead{Proposal 3 --- What Landis Throws In --- Gifts Past the Menu}
  \emph{Your margin: what kind of thing could Landis give when Slack is due?}
  Slack already waives price; the `something extra' should be a thing
  money/price can't buy --- strange, true, and never named in advance. A
  keeper's stock to draw from:
  \begin{itemize}[leftmargin=6mm,itemsep=1mm,topsep=1mm]
    \item \emph{A name not yet earned} --- the true name of a person, ship, or
      thing the crew hasn't met.
    \item \emph{A glimpse down the tide} --- one honest image of a Work's next
      squall before it breaks.
    \item \emph{A ward} --- a one-time `this once, the water lies for you': a
      single future price refunded.
    \item \emph{A door} --- the location of still water somewhere it shouldn't
      be (behind enemy lines, aboard a ship).
    \item \emph{A dead man's sentence} --- the last true thing someone said
      before the crew could ask.
  \end{itemize}
  \vspace{0.5mm}
  Rule of thumb: the gift is \emph{knowledge or leverage}, never raw Power ---
  Landis deals in truth, not drones.
\end{proposalbox}

% =====================================================================
% BOOK III -- WORKS & THE TIDE
% =====================================================================
\clearpage
\booktitle{Council · A Player's Handbook}{WORKS \& THE TIDE}{Book the Third}%
  {how the world keeps moving even while the table sleeps}

\begin{pullquote}A scene is what the crew does with their hands. A Work is what
they set in motion and must then wait upon. Everything in this book is bought
with the one coin the game keeps honest for everyone at once --- time.\end{pullquote}

\sect{3.1 Declaring a Work}
A Work is any labour too large for a single scene --- a bomb built, a treaty
struck, a ruin excavated, a wound healed. It is written on a card the whole
table can see, and it carries three lines and no more:
\begin{itemize}[leftmargin=6mm,label=\textcolor{gold}{\ding{119}},itemsep=1mm,topsep=1mm]
  \item \textbf{THE AIM} --- the plain thing wanted, stated so plainly that its
    finishing cannot be argued.
  \item \textbf{THE HANDS} --- who does it: a person, a department, a people.
  \item \textbf{THE STAKES} --- what is pledged from the Ledger to make it possible.
\end{itemize}

\sect{3.2 The Tide --- Watches}
Time in Council comes in like a tide, on its own, needing no one's permission.
\cut{Every session, every open Work advances one watch --- whether the crew
touched it or not.} \add{A watch is a unit of fictional time, not a unit of
table time. When the keeper turns the watch, every open Work advances one ---
whether the crew touched it or not. A session may pass one watch, or several; a
night's sleep passes a watch, and a long journey may pass more.} Spend a whole
scene labouring at a Work, and it takes \emph{one watch more}. The counts are
public.

\begin{tcbraster}[raster columns=3,raster equal height=rows,raster column skip=3mm,
    sharp corners,boxrule=1pt,colback=white,colframe=navy,coltext=black,
    left=3.5mm,right=3.5mm,top=3mm,bottom=3mm]
  \tcbitem \cardhead{Small Work}{} 1 watch
  \tcbitem \cardhead{A Work}{} 3 watches
  \tcbitem \cardhead{Great Work}{} 5 watches
\end{tcbraster}

\begin{proposalbox}
  \prophead{Proposal 1 --- Turning the Watch --- Keeper Guidance}
  You decided a watch is fictional time, not table time. Open question from your
  margin (p.12): \emph{how many watches does a session pass?} My recommendation
  --- don't hard-quantify it; give the keeper three anchors and let fiction do
  the rest:
  \begin{itemize}[leftmargin=6mm,itemsep=1mm,topsep=1mm]
    \item \emph{Sleep = 1 watch.} A night's rest always turns one.
    \item \emph{Travel by distance.} A hop across the city = 0. A trip off-world
      and back = 1. A long expedition = 2--3, keeper's call.
    \item \emph{A full scene of labour = +1} to the Work worked (already in the
      rules).
  \end{itemize}
  \vspace{0.5mm}
  So a quiet session might pass one watch; a session that sleeps, travels, and
  labours might pass three. \emph{Your boxed note} (``multiple Works per
  session, keeper's discretion'') is already true under this: turning a watch
  advances \emph{every} open Work at once. The thing to decide is whether
  players can ever \emph{force the keeper to turn a watch} --- I'd say no; only
  the fiction turns it, and Forcing (\S3.5) is the one player lever on pace.
\end{proposalbox}

\sect{3.3 The Ledger --- Stakes}
A Work spends more than time; it spends the crew's finite holdings, pledged
under three headings:

\begin{tcbraster}[raster columns=3,raster equal height=rows,raster column skip=3mm,
    sharp corners,boxrule=1pt,colback=white,colframe=navy,coltext=black,
    left=3.5mm,right=3.5mm,top=3mm,bottom=3mm]
  \tcbitem \cardhead{Hands}{} people \& their labour
  \tcbitem \cardhead{Power}{} resources, the means
  \tcbitem \cardhead{Standing}{} goodwill with peoples
\end{tcbraster}
\emph{A stake cannot be pledged twice} --- and \textbf{to wound a pledged stake
is to stall everything pledged to it.} This is how a single blow lands
everywhere.

\begin{proposalbox}
  \prophead{Proposal 2 --- Giving the Ledger Meat --- Hard Hands \& Power, Hidden Standing}
  Your notes (p.2 ``needs meat,'' p.11 ``quantified \dots\ city resource
  management'') and your answer both point the same way. Proposed split:
  \begin{itemize}[leftmargin=6mm,itemsep=1mm,topsep=1mm]
    \item \emph{Hands \& Power are counted, and public.} Give the expedition a
      pool --- e.g. Hands as named departments/teams, Power as concrete stores
      (ZPM charge, C4, jumpers, drones). A Work pledges specific units; the crew
      can always see what's committed and what's free. This is the `city
      resource management' layer.
    \item \emph{Standing is not a number the players see.} Per your answer: the
      keeper tracks each people's regard (Genii, Hoffan, Athosian\dots) on a
      hidden scale; players must \emph{intuit} it from how factions behave. The
      keeper knows whether a faction will lift a hand; the crew has to read the
      room.
  \end{itemize}
  \vspace{0.5mm}
  Open sub-question I'd flag: do you want Standing tracked per-faction as a small
  hidden dial (e.g. $-2\dots+2$), or purely as keeper prose? A dial makes ``will
  they help?'' rulings consistent; prose keeps it loose. I lean \emph{hidden
  dial, prose to the players.}
\end{proposalbox}

\sect{3.4 Squalls}
\cut{Every Great Work} \add{Every Great Work --- and, at the keeper's
discretion, some lesser Works ---} is declared with \emph{one sealed
complication} --- its squall --- set face-down beside the card. It breaks at the
\emph{halfway watch}, or sooner if the weave calls for it.

\sect{3.5 Forcing a Work}
Once per session, a crew may drive a single Work \emph{one extra watch} by
paying one of two coins:

\begin{tcbraster}[raster columns=2,raster equal height=rows,raster column skip=4mm,
    sharp corners,boxrule=0pt,colback=cream,colframe=cream,coltext=black,
    left=3.5mm,right=3.5mm,top=3mm,bottom=3mm]
  \tcbitem \textbf{A TRUTH} --- confessed aloud, canon the instant it is spoken.
    \add{Or a truth made true \emph{of the Council} --- a new law passed.}
  \tcbitem \textbf{A FAVOR OWED} --- sealed, called due whenever the keeper pleases.
\end{tcbraster}
Note the mark of haste: a forced Work \cut{always} \add{still} delivers its squall.

\begin{proposalbox}
  \prophead{Proposal 4 --- ``A Favor Owed --- Sealed.'' What is Sealed?}
  You struck `sealed' and asked \emph{what is there to seal?} Two clean readings
  --- pick one:
  \begin{itemize}[leftmargin=6mm,itemsep=1mm,topsep=1mm]
    \item \emph{(a) Sealed = undefined-until-called.} The favor has no named
      price now; the keeper writes it face-down and reveals it whenever they
      call it due. What's `sealed' is the \emph{cost}. (Keeps the dread; matches
      Squalls.)
    \item \emph{(b) Drop the seal.} A favor owed is just an IOU the keeper can
      call anytime --- no envelope, no mystery. Simpler at the table.
  \end{itemize}
  \vspace{0.5mm}
  I lean \emph{(a)}: it makes `a favor owed' scarier than `a truth' and gives
  you a sealed prop to hand over, which fits Council's whole reveal aesthetic.
\end{proposalbox}

\sect{3.6 Completion \& Work Orders}
A finished Work stops being a card and becomes a \emph{plain fact of the world.}
Between sessions the crew may send \emph{Work Orders} --- written instructions
to the hands who do the labour --- and the keeper answers in those hands' own
voices next the table sits.

\begin{proposalbox}
  \prophead{Proposal 5 --- Between-Session Play --- Orders That Resolve Off-Table}
  Your green note (p.13): \emph{a way for things to progress between sessions
  without wrecking the timeline, and encouraging between-session play.} Work
  Orders already point here; proposed expansion:
  \begin{itemize}[leftmargin=6mm,itemsep=1mm,topsep=1mm]
    \item Between sessions, any player may file \emph{Standing Orders} to the
      hands they command --- written instructions (build, scout, treat, dig).
    \item These \emph{don't turn extra watches} (the timeline is safe): they
      resolve as of the \emph{next} watch the fiction turns, and the keeper
      answers in the hands' own voice (letterhead) at session start.
    \item A player who files orders may open the session already mid-scene on
      their result --- the reward for playing between sessions.
  \end{itemize}
  \vspace{0.5mm}
  This is a natural fit for the Council website: orders in, letterhead replies out.
\end{proposalbox}

\sect{3.7 The Divine Work --- The Great Tier}
Above the Great Work stands one tier more, reserved for labours that span a
whole telling\add{, or a season}. A \emph{Divine Work} is not a card the crew
declares; it is a weather the keeper sets moving beneath everything:

\begin{tcbraster}[raster columns=1,raster row skip=2mm,sharp corners,boxrule=0pt,
    colback=navy,colframe=navy,coltext=cream,left=4mm,right=4mm,top=2.5mm,bottom=2.5mm,
    before skip=2mm,after skip=2mm]
  \tcbitem \textbf{No fixed count.} No watch-total set at declaring, no
    countdown shown \add{--- even if one is quietly kept.}
  \tcbitem \textbf{Many squalls, not one.} Several sealed \emph{beats}, opened
    one at a time.
  \tcbitem \textbf{A beat may rewrite the card.} A beat can recast the aim, add
    a stake, release another.
  \tcbitem \textbf{Stakes bind and loose over time.} Hands and Standing attach
    and detach across its length.
\end{tcbraster}
A Divine Work ends not on a watch but on a \emph{turning}. Even here the one law
holds: \emph{mercy}.

\begin{notebox}
  {\sffamily\bfseries\footnotesize\textcolor{gold!75!black}%
    {\textls[100]{HOW A STORM SPEAKS --- A STANDING HOUSE STYLE}}}\par\vspace{1.5mm}
  Every squall, every beat, every hard turn is delivered as an \emph{in-fiction
  report on the letterhead of whoever bore it.} Never an omniscient voice from
  the clouds. In Council, the bad news always arrives signed
  \add{--- witnessed, not narrated. Testimony is king.}
\end{notebox}

\begin{proposalbox}
  \prophead{Proposal 6 --- Trigger Points \& Non-Step-able Countdowns --- The Divine Work Engine}
  You don't have the shape yet and asked for options. Here's a model that fits
  `The Storm shouldn't tick until the Genii are ready to be enemies.'

  \emph{The idea.} A Divine Work is a chain of \emph{beats}. A beat does nothing
  until its \emph{Trigger} is met. When it fires, that beat's \emph{watch-clock}
  starts counting --- and only then does time bite.

  \emph{A beat has three lines: Trigger} (the fictional condition that arms it),
  \emph{Clock} (how many watches once armed), \emph{Break} (what the beat does
  --- usually a signed squall, sometimes a rewrite of the card).

  \emph{`Non-step-able' means two guarantees:} (1) an un-armed beat can't be
  advanced, forced, or grace'd forward --- no lever moves a clock that hasn't
  started; (2) once armed, the clock is the keeper's, not the tide's default ---
  it ticks on its own terms.

  \emph{Storm, worked example:} Beat I \emph{Outer Bands} --- Trigger: Genii
  gain a foothold · Clock: none (mood only). Beat III \emph{Eye} --- Trigger:
  storm makes landfall on Atlantis · Clock: 3 watches · Break: shield fails,
  signed by Radek. Beat IV --- Trigger: Beat III resolves · Break: Earth's ship
  arrives with McKay + partial ZPM.

  \emph{Decision I need:} is a beat's clock \emph{visible} to players once armed
  (a ticking dread they can race) or \emph{hidden} (they feel the squeeze but
  don't see the number)? I'd make it the keeper's choice per beat.
\end{proposalbox}

\begin{proposalbox}
  \prophead{Proposal 7 --- Three Tiers of Divine Work --- Keeper Guidance (Not Hard Tiers)}
  You want these as guidance, not a rigid parallel to Small/Work/Great. Proposed
  as a keeper sidebar:
  \begin{itemize}[leftmargin=6mm,itemsep=1mm,topsep=1mm]
    \item \emph{An Episode} --- one session's worth; a single strong beat.
      \emph{Ex: the dying ZPM on the children's planet.}
    \item \emph{An Arc} --- spans several sessions; a handful of chained beats.
      \emph{Ex: The Storm; poisoning the well; the first Wraith siege.}
    \item \emph{A Saga} --- one or more seasons; beats that sleep for a long
      time between arming. \emph{Ex: the Wraith threat; Taxe's Federation;
      powering the city and getting home.}
  \end{itemize}
  \vspace{0.5mm}
  Same engine (beats + triggers) at every tier --- only the span and the
  patience differ.
\end{proposalbox}

\begin{proposalbox}
  \prophead{Proposal 8 --- ``Dependent Visible Works'' --- What I Meant, and My Recommendation}
  This was my jargon; here's the plain version. A Divine Work is \emph{weather}
  the players can't grab directly. A \emph{dependent visible Work} is an
  ordinary Work card the storm throws off that they \emph{can} grab --- the
  handhold. The Storm (Divine) can't be `solved'; but \emph{Shield's Heart},
  \emph{Project Hammer}, and \emph{The Catch} are visible Works the crew can
  pour Hands into to change how the storm lands.

  \emph{My recommendation, not a hard rule:} every armed beat should expose
  \emph{at least one} visible Work the players can act on --- otherwise the beat
  is something that happens \emph{to} them, which is railroading by another
  name. Make it a keeper habit, not a law.
\end{proposalbox}

% =====================================================================
% BOOK IV -- GOVERNANCE
% =====================================================================
\clearpage
\booktitle{Council · A Player's Handbook}{GOVERNANCE}{Book the Fourth}%
  {who may command a thing to happen --- and what it costs to choose}

\begin{pullquote}The game is named for this book. A Council is not a throne ---
it is a table, where authority is shared, choices are witnessed, and the cost of
command is that everyone remembers what you decided.\end{pullquote}

\sect{4.1 The Council}
A Council is a body of shared command \hlt{with a single seat} that holds the
\emph{final say}. Its virtue is not efficiency --- it is \emph{the record.}
Every vote and every dissent is entered and remembered; a voice overruled is not
silenced but \emph{minuted}.

\begin{pullquote}A king may be wrong in private. A chair may only be wrong in
front of the table.\end{pullquote}

\begin{proposalbox}
  \prophead{Proposal 9 --- The Fasces --- An Optional Name \& Rule for the Final Seat}
  Your p.15 note: call the single-seat holder the one who \emph{holds the
  Fasces}, rotating each session, and --- like the Roman consuls --- the others
  in rotation may \emph{veto}. Proposed optional rule:
  \begin{itemize}[leftmargin=6mm,itemsep=1mm,topsep=1mm]
    \item The final-say seat rotates each session among those eligible.
    \item Any other rotation-holder may spend a \emph{veto} to block a single
      ruling (once per session), forcing a Council vote instead.
  \end{itemize}
  \vspace{0.5mm}
  \emph{Your wide/narrow question:} yes --- at a \emph{wide table} the keeper can
  absolutely take the Fasces on some sessions (an in-fiction chair, not the GM
  hat), which is a lovely way to put the keeper's avatar on the record. At a
  \emph{narrow table}, rotation is usually just among the player-principals and
  the keeper stays the world. Tie it to the \S4.2 dial.
\end{proposalbox}

\begin{proposalbox}
  \prophead{Proposal 10 --- The Minute --- Record-Keeping (and Your Voice-Memo Question)}
  You circled `on the record' (\S4.5) and asked if minuting is practical ---
  \emph{can we use voice memos?} Proposed light system, \emph{The Minute}: after
  any vote, override, or Directive, one line goes in the log --- \emph{who
  decided, what, who dissented}. That's all `minuted' needs to mean. A phone
  voice-memo is a perfectly good Minute.

  \emph{On tooling --- straight answer:} I can't transcribe an audio file
  directly in this workspace (no speech-to-text here). But the loop you want
  works fine: run tonight's recording through any transcription app, upload the
  \emph{transcript} (txt/doc), and I'll keep a running game log \emph{and} a
  formal Council Minutes ledger from it --- votes, Directives, Works advanced,
  squalls broken. If you'd rather, dictate a recap (like you just did) and I'll
  minute that. I'd suggest we start a \texttt{Council\_Minutes.md} in the
  project and append after each session.
\end{proposalbox}

\sect{4.2 Who Sits --- A Dial Set at Setup}
Council does not decree who holds a seat; the table sets that dial when it seats
its own Council:

\begin{tcbraster}[raster columns=2,raster equal height=rows,raster column skip=4mm,
    sharp corners,boxrule=0pt,colback=navy,coltext=cream,
    left=3.5mm,right=3.5mm,top=3mm,bottom=3mm]
  \tcbitem[colframe=gold,toprule=1.4pt] \cardhead{The Wide Table}{} Every
    standing head holds a seat and a vote. Loud, slow, truly divided. Choose
    this for intrigue.
  \tcbitem[colframe=teal,toprule=1.4pt] \cardhead{The Narrow Table · Default}{}
    Only the few at the heart of the story sit and vote; the heads report and
    advise. Choose this for pace --- and by default, \emph{this is where a
    Council begins.}
\end{tcbraster}
Council seats a narrow table by default, and widens only when a table decides
the intrigue of a true vote is worth the pace it costs.

\sect{4.3 The Second Chair --- Deputies}
A deputy may \hlt{hold a Work at its present watch}, but may not \emph{advance}
it. To move a Work forward wants a Work Order or a Directive from the seat of
command.

\begin{proposalbox}
  \prophead{Proposal 11 --- Reversion --- What Happens to an Unheld Work}
  Your p.15 question: \emph{without a deputy, does a pulled Work revert one step
  each watch?} I think your instinct is the right rule, and it's better than a
  flat `freeze' --- it gives interrupts real teeth. Proposed:
  \begin{itemize}[leftmargin=6mm,itemsep=1mm,topsep=1mm]
    \item A Work whose head is \emph{pulled} (\S4.4) and left \emph{unheld slips
      back one watch} each watch that passes --- the labour decays without its
      hand.
    \item A \emph{deputy} holding it stops the slip (freeze at present watch)
      but can't advance it.
    \item A \emph{Directive} (\S4.5) keeps the head on task, so it neither slips
      nor pauses.
  \end{itemize}
  \vspace{0.5mm}
  This answers your `what about Carson pinned on the Hoffan cure?' --- pull
  Carson to an emergency and, unless Keely (deputy) holds it, the cure work rots
  a watch at a time until he's back or you Directive him to stay. That's a real
  cost, which is the point.
\end{proposalbox}

\sect{4.4 Interrupts}
An emergency that names a head \emph{pulls them off their Work automatically.}
The card \emph{freezes} until released, or a deputy steps in to \hlt{hold it}
\add{--- keeping it from slipping back}.

\sect{4.5 The Override}
Only the seat of final say may keep a pulled head at their task regardless, by
issuing a \emph{Directive} --- a \emph{truth committed aloud, on the record.}

\sect{4.6 Absence Beyond Reach}
When a head is truly beyond word, the need is met anyway, by whoever stands
present, \emph{without that head's insight} --- it shows as something
\emph{concrete and specific}, never a clean sealed squall.

\begin{proposalbox}
  \prophead{Proposal 12 --- \S4.6 Absence Beyond Reach --- Proposed Rewrite}
  You flagged this as thin. It's the twin of Interrupts: Interrupts pull a head
  who's \emph{here}; Absence covers a head who is simply \emph{unreachable}
  (off-world, ascended, missing). Proposed replacement text:
  \begin{tcolorbox}[blanker,borderline west={2pt}{0pt}{teal},left=4mm,
      top=1mm,bottom=1mm,before skip=2mm,after skip=2mm]
    \itshape When a head is beyond all word --- off-world, ascended, or lost ---
    the table cannot wait on their insight, and the need is met by whoever
    stands present. The answer still comes, but it comes rough: concrete and
    specific, shaped by the wrong hands, never the clean sealed squall the
    absent expert would have drawn. A botany question answered by a soldier gets
    a soldier's answer. Mark such Works: when the head returns, they may find the
    labour done crooked, and worth re-doing.
  \end{tcolorbox}
  The distinction to keep: Interrupt = freeze/slip (they're here, just busy);
  Absence = it happens anyway, badly (they're gone).
\end{proposalbox}

% =====================================================================
% BOOK V -- TEACH & SHIP
% =====================================================================
\clearpage
\booktitle{Council · A Player's Handbook}{TEACH \& SHIP}{Book the Fifth}%
  {for the keeper --- how to raise a Council game from nothing, and run it well}

\begin{pullquote}The first four books are the game. This one is the craft of it
--- how to stand a table up in an evening, and the three small arts that make a
keeper of you: writing Seals, giving grace, and wounding without cruelty.\end{pullquote}

\sect{5.1 Session Zero --- Raising a Game in One Sitting}
Council boots in a single evening. Five moves, in order:
\begin{enumerate}[leftmargin=8mm,itemsep=1.3mm]
  \item \emph{Rig the world.} Agree the fiction --- the setting, its peoples,
    what Power and Standing are made of here.
  \item \emph{Seat the Council.} Decide who holds the final say, and set the
    dial of \S4.2 --- \emph{narrow by default}, or wide for intrigue.
  \item \emph{Lay the first Ledger.} Write two or three opening Works.
  \item \emph{Teach the Water.} Tell the crew, plainly, of Landis.
  \item \emph{Promise hope.} Say aloud the one creed. Then begin.
\end{enumerate}

\begin{proposalbox}
  \prophead{Proposal 13 --- The One Creed --- A Poetic Draft, \& the Session-Zero Ledger Step}
  Two things land here. First, your `write a poetic creed' --- a first swing, to
  be said aloud at Session Zero (\S5.1, move 5):
  \begin{tcolorbox}[blanker,before skip=2mm,after skip=2mm]
    \centering\itshape
    We keep no dice, and we keep no fear.\\
    What we cannot answer, we will pay for true;\\
    what we cannot pay, we will be given.\\
    No door is locked that hope can't open ---\\
    and the only door that locks from the inside is despair.\\
    So we strive, and we pay, or we are given ---\\
    and we keep hope, whatever the weather.
  \end{tcolorbox}
  Second, move 3 (`Lay the first Ledger') needs expanding once P2 lands --- it
  should walk the keeper through setting starting Hands/Power pools and, if you
  want them, \emph{research tiers} (what the expedition can build now vs. what
  needs a Work first). Flagging; I'll draft it after you rule on P2.
\end{proposalbox}

\sect{5.2 The First Art --- Writing Seals}
For any obstacle worth being stuck on, \emph{write three clues before play and
seal them}: the Nudge, the Bearing, the Answer. Write them \emph{before} ---
never under the pressure of a stuck table, when your judgment is worst.

\begin{proposalbox}
  \prophead{Proposal 14 --- Sourcing Clues --- and Which Letterhead They Ride In On}
  Your p.17 note: clues can come from many sources --- prep them too, at least
  which source and which letterhead. Proposed addition to \S5.2:
  \begin{itemize}[leftmargin=6mm,itemsep=1mm,topsep=1mm]
    \item When you seal a clue, also note its \emph{source} --- who knows it,
      and how it would reach the crew.
    \item Give it a \emph{letterhead}: a clue from Beckett arrives as a medical
      note; from the Genii, as an intercept; from Landis, as a line in still
      water. Same clue, different voice, different trust.
  \end{itemize}
  \vspace{0.5mm}
  This dovetails with the Signed Squall (\S5.4): in Council, \emph{everything}
  that arrives --- good news or bad --- arrives signed.
\end{proposalbox}

\sect{5.3 The Second Art --- Giving Grace}
Give it \emph{quietly}. Give it toward the \emph{stuck and the hopeful.} Do not
tally it, do not let it be demanded. Its one hard law is that a truly aground
crew always receives it.

\sect{5.4 The Third Art --- The Signed Squall}
Deliver bad news as a \emph{report on the letterhead of whoever bore it.} In
Council, bad news always arrives signed.

\sect{5.5 On Hope \& Harm --- The Tone of the Game}
Council is built to reward hope. The game \emph{may wound} --- it may cost
dearly, and grieve truly. But it breaks only the one who has \emph{stopped
hoping.}

\begin{pullquote}Despair is the only door that locks from the inside.\end{pullquote}

\sect{5.6 What to Print, What to Keep}
Two bundles make a Council game ready for a table:

\begin{tcbraster}[raster columns=2,raster equal height=rows,raster column skip=4mm,
    sharp corners,boxrule=0pt,left=3.5mm,right=3.5mm,top=3mm,bottom=3mm]
  \tcbitem[colback=cream,colframe=teal,coltext=black,toprule=1.4pt]
    \cardhead{}{The Crew's Packet} This handbook, and Landis's two handouts. All
    a player needs. \add{(For the Council of Un: the crew's Standing Orders for
    the expedition ride here too.)}
  \tcbitem[colback=navy,colframe=gold,coltext=cream,toprule=1.4pt]
    \cardhead{}{The Keeper's Papers} Sealed Seals, hidden stakes, Divine Work
    beats. Shown to no one until its watch turns.
\end{tcbraster}

\begin{tcolorbox}[enhanced,sharp corners,colback=navy,colframe=navy,
    coltext=cream,boxrule=0pt,borderline west={2.5pt}{0pt}{gold},
    left=5mm,right=5mm,top=3.5mm,bottom=3.5mm]
  \textbf{\textcolor{gold}{The Living Ledger}} --- a place to run Works between
  sessions: on paper, or (better) online at the Council website.
\end{tcolorbox}

\begin{center}\itshape Print the first. Keep the second. \add{Run the third.}\end{center}

\vfill
\begin{center}\deco\end{center}

\end{document}
